
\begin{document}
\extrawidth{0.5in} \extrafootheight{-0in} \pagestyle{headandfoot}
\headrule \header{\textbf{CS2311 - Spring 2020}}{\textbf{HW
 6 \ifprintanswers - Solutions \fi}}{\textbf{Due: Sun. 03/08/20}} \footrule \footer{}{Page \thepage\
of \numpages}{}

\ifprintanswers
\noindent \textbf{Instructions:} All assignments are due \underline{by \textbf{midnight} on the due date} specified.  Assignments should be typed and submitted as a PDF.  Every student must write up their own solutions in their own manner.

\medskip
\noindent You should \underline{complete all problems}, but \underline{only a subset will be graded} (which will be graded is not known to you ahead of time). 
\else
\noindent \textbf{Instructions:} All assignments are due \underline{by \textbf{midnight} on the due date} specified.  Every student must write up their own solutions in their own manner.

\noindent Please present your solutions in a clean, understandable
manner.  Use the provided files that give mathematical notation in Word, Open Office, Google Docs, and \LaTeX. 

\noindent Assignments should be typed and submitted as a PDF.   

\noindent You should \underline{complete all problems}, but \underline{only a subset will be graded} (which will be graded is not known to you ahead of time). 
\fi

\begin{questions}


\section*{Graphs}

\ugquestion{3} Rosen Ch 10.5 \#56, p. 742.
  \ifprintanswers
        \vspace{-10pt}
  \fi
  \begin{solution}
    \includegraphics[width=1in]{figs/rosen_10_5_56}
  \end{solution}


  \ifprintanswers
        \vspace{-10pt}
  \fi
\gquestion{2}{2}{all} Rosen Ch 10.5 \#58(a), p. 742.
  \ifprintanswers
        \vspace{-10pt}
  \fi
  \begin{solution}
  \textit{A \textbf{knight’s tour} is a sequence of legal moves by a knight starting at some square and visiting each square exactly once. We can model knights tours using the graph that has a vertex for each square on the board, with an edge connecting two vertices if a knight can legally move between the squares represented by these vertices.}

  \begin{enumerate}[label=(\alph*),itemsep=0pt,parsep=0pt,topsep=0pt,partopsep=0pt]
    \item \textit{Show that finding a knight’s tour on an $m \times n$ chessboard is equivalent to finding a Hamilton path.} \\
    Let the graph represent legal moves of the knights, where the vertices are the squared on the board and the edges between those squares where moves are legal. 

    By definition, a knight’s tour visits each square exactly once, therefore traverses a simple path that includes each vertex once; which fits the definition of a Hamilton path. 

    % \item \textit{Show that finding a reentrant knight’s tour on an m × n chessboard is equivalent to finding a Hamilton circuit on the corresponding graph.}  \\

    \end{enumerate}

  \end{solution}



  \ifprintanswers
        \vspace{-10pt}
  \fi
\gquestion{2}{2}{all} Rosen Ch 10.7 \#14, p. 761.
  \ifprintanswers
        \vspace{-10pt}
  \fi
  \begin{solution}
    You are given $e = 30$ and $r = 20$. Therefore, rewriting Euler's formula: 
      $r = e - v + 2$ so $v = e - r + 2 = 30 - (20) + 2 = 12$
  \end{solution}



  \ifprintanswers
        \vspace{-10pt}
  \fi
\ugquestion{2} Rosen Ch 10 Supplementary Exercises \#19, p. 774.
  \ifprintanswers
        \vspace{-10pt}
  \fi
  \begin{solution}
    The graph does not contain $K_5$ or any larger complete graphs.  It has the following cliques: 
    \begin{itemize}[itemsep=0pt,parsep=0pt,topsep=0pt,partopsep=0pt]
      \item $K_4 - \{b, c, e, f\}$ 
      \item $K_3 - \{a, b, g\}$, $\{a, d, g \}$, $\{d, e, g\}$, $\{b, e, g\}$
    \end{itemize}
  \end{solution}



  \ifprintanswers
        \vspace{-10pt}
  \fi
\gquestion{2}{2}{all} Rosen Ch 10 Supplementary Exercises \#20, p. 774.
  \ifprintanswers
        \vspace{-10pt}
  \fi
  \begin{solution}
  The graph does not contain $K_6$ or any larger complete graphs. It has the following cliques: 
  \begin{itemize}[itemsep=0pt,parsep=0pt,topsep=0pt,partopsep=0pt]
    \item $K_5 -  \{c, e, g, h, i\}$
    \item $K_4 -  \{a, b, c, e\}$, $\{c, d, e, g\}$
    \item $K_3$ - $\{a,e,f\}, \{e, f, g\}$
  \end{itemize}
  \end{solution}


  \ifprintanswers
        \vspace{-10pt}
  \fi
\gquestion{2}{2}{all} Rosen Ch 10 Supplementary Exercises \#22, p. 775.
  \ifprintanswers
        \vspace{-10pt}
  \fi
  \begin{solution}
    The unique minimum dominating set is $\{e \}$
  \end{solution}



\section*{Logic}
\ifprintanswers
\vspace{-10pt}
\uplevel{\color{blue} \large For each question, there are many possible solutions.  I have typically provided one or two common solutions to these problems. }
\fi


\gquestion{34}{24}{a-c} Prove the following statement is a tautology without using truth tables (use the logical equivalences from Table 6-8 of the book).  Justify each step with the law used.  Model the solutions in the style of Examples 6-8, pp. 29-30 of the book.


\begin{enumerate}[label=(\alph*),itemsep=0pt,parsep=0pt,
topsep=0pt,partopsep=0pt]
    \item (6 points) $p \rightarrow (\neg q \rightarrow r) \equiv \neg(q \vee r)
    \rightarrow \neg p$

    % f15, f14
    \item (8 points) $p \ra (q \vee r) \equiv (p \wedge \neg q) \ra r$

    % f14
    \item (10 points) $ [ \neg p \wedge (p \vee q) ] \ra q \equiv \mathbf{T}$
    
    % f16
    \item (10 points) $[(p \vee q) \wedge (p \ra r) \wedge (q \ra r)] \ra r \equiv \mathbf{T}$
\end{enumerate}
\begin{solution}
  \begin{enumerate}[label=(\alph*),itemsep=0pt,parsep=0pt,
topsep=0pt,partopsep=0pt]
    \item $p \rightarrow (\neg q \rightarrow r) \equiv 
      \neg(q \vee r) \rightarrow \neg p$
    \begin{align*}
      p & \rightarrow (\neg q \rightarrow r) \\
      & \equiv p \ra (\neg \neg q \vee r) \tag{Table 7, rule 1} \\
      & \equiv p \ra (q \vee r) \tag{Double Negation} \\
      & \equiv \neg (q \vee r) \ra \neg p \tag{Table 7, rule 2} \\
    \end{align*}

    \vspace{-25pt}
    \begin{center}
    OR
    \end{center}

    \vspace{-25pt}
    \begin{align*}
      p & \rightarrow (\neg q \rightarrow r) \\
       & \equiv p \ra (q \vee r) \tag{Table 7, rule 3} \\
       & \equiv \neg (q \vee r) \ra \neg p  \tag{Table 7, rule 2}
    \end{align*}

    \item $p \ra (q \vee r) \equiv (p \wedge \neg q) \ra r$
    \begin{align*}
        p &\ra (q \vee r) \equiv (p \wedge \neg q) \ra r \\
            & \equiv \neg p \vee (q \vee r) \tag{Table 7, rule 1} \\
            & \equiv (\neg p \vee q) \vee r) \tag{Associative} \\
            & \equiv \neg (\neg p \vee q) \ra r \tag{Table 7, rule 1} \\
            & \equiv (\neg \neg p \wedge \neg q) \ra r \tag{DeMorgan's} \\
            & \equiv (p \wedge \neg q) \ra r \tag{Double Negation} 
    \end{align*}
    
    \begin{center}
    OR
    \end{center}

    \begin{align*}
        p &\ra (q \vee r) \equiv (p \wedge \neg q) \ra r \\
        & \equiv (p \ra q) \vee (p \ra r) \tag{Table 7, rule 8} \\
        & \equiv (\neg p \vee q) \vee (p \ra r) \tag{Table 7, rule 1} \\
        & \equiv (\neg p \vee q) \vee (\neg p \vee r) \tag{Table 7, rule 1} \\
        & \equiv \neg p \vee q \vee \neg p \vee r \tag{drop parenthesis} \\
        & \equiv \neg p \vee \neg p \vee q \vee r \tag{Commutative} \\
        & \equiv \neg p \vee q \vee r \tag{Idempotent Law} \\
        & \equiv (\neg p \vee q) \vee r \tag{add parenthesis} \\
        & \equiv \neg \neg (\neg p \vee q) \vee r \tag{Double Negation} \\
        & \equiv \neg (\neg \neg p \wedge \neg q) \vee r \tag{DeMorgans} \\
        & \equiv \neg (p \wedge \neg q) \vee r \tag{Doulbe Negation} \\
        & \equiv (p \wedge \neg q) \ra r \tag{Table 7, rule 1}
    \end{align*}

    \item $ [ \neg p \wedge (p \vee q) ] \ra q \equiv \mathbf{T}$
    % There are many possible solutions; two are given.
    \small
    \begin{align*}
        [ \neg p & \wedge (p \vee q) ] \ra q \\
        & \equiv \neg [ \neg p \wedge (p \vee q) ] \vee q \tag{Table 7, rule 1} \\
        & \equiv \neg \neg p \vee \neg (p \vee q) \vee q \tag{DeMorgan's law} \\
        & \equiv p \vee \neg (p \vee q) \vee q \tag{Double Negation} \\
        & \equiv p \vee q \vee \neg (p \vee q) \tag{Commutative} \\
        & \equiv (p \vee q) \vee \neg (p \vee q) \tag{add in parenthesis}\\
        & \equiv \mathbf{T} \tag{Negation} 
    \end{align*}

    \begin{center}
    OR
    \end{center}

    \vspace{-25pt}
    \begin{align*}
        [ \neg p & \wedge (p \vee q) ] \ra q \\
        & \equiv \neg [ \neg p \wedge (p \vee q) ] \vee q \tag{Table 7, rule 1} \\
        & \equiv \neg \neg p \vee \neg (p \vee q) \vee q \tag{DeMorgan's law} \\
        & \equiv p \vee \neg (p \vee q) \vee q \tag{Double Negation} \\
        & \equiv p \vee q \vee \neg (p \vee q) \tag{Commutative} \\
        & \equiv p \vee q \vee (\neg p \wedge \neg q) \tag{DeMorgan's} \\
        & \equiv p \vee (q \vee \neg p) \wedge (q \vee \neg q) \tag{Distributive} \\
        & \equiv p \vee (q \vee \neg p) \wedge \mathbf{T} \tag{Negation} \\
        & \equiv p \vee (q \vee \neg p) \tag{Identity} \\
        & \equiv p \vee (\neg p \vee q) \tag{Commutative} \\
        & \equiv (p \vee \neg p) \vee q \tag{Associative} \\
        & \equiv \mathbf{T} \vee q \tag{Negation} \\
        & \equiv \mathbf{T} \tag{Domination}
    \end{align*}
    

    \item $[(p \vee q) \wedge (p \ra r) \wedge (q \ra r)] \ra r \equiv \mathbf{T}$ 
    \begin{align*}
      [(p \vee q) &\wedge (p \ra r) \wedge (q \ra r)] \ra r \\
      &\equiv \neg[(p \vee q) \wedge (\neg p \vee r) \wedge (\neg q \vee r)] \vee r \tag{Table 7, rule 1, x3} \\
      &\equiv \neg(p \vee q) \vee \neg(\neg p \vee r) \vee \neg(\neg q \vee r) \vee r \tag{DeMorgan's} \\
      &\equiv (\neg p \wedge \neg q) \vee (\neg\neg p \wedge \neg r) \vee (\neg\neg q \wedge \neg r) \vee r \tag{DeMorgan's, x3} \\
      &\equiv (\neg p \wedge \neg q) \vee (p \wedge \neg r) \vee (q \wedge \neg r) \vee r \tag{Double Negation, x2} \\
      &\equiv (\neg p \wedge \neg q) \vee (p \wedge \neg r) \vee r \vee (q \wedge \neg r) \tag{Commutative Law} \\
      &\equiv (\neg p \wedge \neg q) \vee (p \wedge \neg r) \vee ((r \vee q) \wedge (r \vee \neg r)) \tag{Distributive Law} \\
      &\equiv (\neg p \wedge \neg q) \vee (p \wedge \neg r) \vee ((r \vee q) \wedge \mathbf{T}) \tag{Negation Law} \\
      &\equiv (\neg p \wedge \neg q) \vee (p \wedge \neg r) \vee r \vee q  \tag{Identity Law} \\
      &\equiv q \vee (\neg p \wedge \neg q) \vee r \vee (p \wedge \neg r) \tag{Commutative Law, x2} \\
      &\equiv ((q \vee \neg p) \wedge (q \vee \neg q)) \vee ((r \vee p) \wedge (r \vee \neg r)) \tag{Commutative Law, x2} \\
      &\equiv ((q \vee \neg p) \wedge \mathbf{T}) \vee ((r \vee p) \wedge \mathbf{T}) \tag{Negation Law, x2} \\
      &\equiv (q \vee \neg p) \vee (r \vee p)  \tag{Identity Law, x2} \\
      &\equiv q \vee \neg p \vee r \vee p  \tag{Associative Law} \\
      &\equiv p \vee  \neg p \vee q \vee r  \tag{Commutative Law} \\
      &\equiv \mathbf{T} \vee q \vee r  \tag{Negation Law} \\
      &\equiv \mathbf{T} \tag{Domination Law} \\
    \end{align*}

    An alternative solution is:
    \begin{align*}
      [(p \vee q) &\wedge (p \ra r) \wedge (q \ra r)] \ra r \\
       &\equiv [(p \vee q) \wedge ((p \vee q) \ra r)] \ra r \tag{Table 7, rule 7} \\
       &\equiv [(p \vee q) \wedge (\neg (p \vee q) \vee r)] \ra r \tag{Table 7, rule 1} \\
       &\equiv [((p \vee q) \wedge \neg (p \vee q)) \vee ((p \vee q) \wedge r)] \ra r \tag{Distributive} \\
       &\equiv [\mathbf{F} \vee ((p \vee q) \wedge r)] \ra r \tag{Negation} \\
       & \equiv [((p \vee q) \wedge r) \vee \mathbf{F}] \ra r \tag{Commutative} \\
       & \equiv [((p \vee q) \wedge r)] \ra r \tag{Identity} \\
       & \equiv [(p \vee q) \wedge r] \ra r \tag{drop extra ()} \\
       & \equiv \neg [(p \vee q) \wedge r] \vee r \tag{Table 7, rule 1} \\
       & \equiv  [\neg (p \vee q) \vee \neg r] \vee r \tag{DeMorgans} \\
       & \equiv  \neg (p \vee q) \vee [\neg r \vee r ] \tag{Associative} \\
       & \equiv \neg (p \vee q) \vee [r \vee \neg r] \tag{Commutative} \\
       & \equiv \neg (p \vee q) \vee \mathbf{T} \tag{Negation} \\
       & \equiv \mathbf{T} \tag{Domination} 
    \end{align*}

  \end{enumerate}
\end{solution}



\gquestion{8}{8}{all}  Show the following expressions are logically equivalent. 
\begin{quote}
  If you are not lazy, then you work hard or you are clever and lazy. \\[3pt]
  You work hard or you are lazy.
\end{quote}

Use the following propositions to represent the concepts in the expression. Let 
\begin{itemize}[itemsep=0pt,parsep=0pt,topsep=0pt,partopsep=0pt]
  \item $l$ be ``you are lazy'' 
  \item $w$ be ``you work hard''
  \item $c$ be ``you are clever''
\end{itemize}

\begin{solution}
First, translate the expressions into logic. 

$$ \neg l \ra (w \vee (c \wedge l)) \equiv w \vee l $$

\begin{align*}
  \neg l \ra & (w \vee (c \wedge l)) \\
    & \equiv l \vee (w \vee (c \wedge l)) \tag{Table 7.3} \\
    & \equiv l \vee ((c \wedge l) \vee w) \tag{Commutative} \\
    & \equiv (l \vee (c \wedge l)) \vee w \tag{Associative} \\
    & \equiv l \vee w \tag{Absorption} \\
    & \equiv w \vee l \tag{Commutative}
\end{align*}

\end{solution}




% \gquestion{8}{8}{all} Prove the following statement is a tautology without using truth tables (use the logical equivalences from Table 6-8 of the book).  Justify each step with the law used.  Model the solutions in the style of Examples 6-8, pp. 29-30 of the book.
% \[ [ \neg p \wedge (p \vee q) ] \ra q \]
%     \ifprintanswers
%         \vspace{-30pt}
%     \fi
% \begin{solution}  There are many possible solutions; two are given.
% \small
%     \begin{align*}
%         [ \neg p & \wedge (p \vee q) ] \ra q \\
%         & \equiv \neg [ \neg p \wedge (p \vee q) ] \vee q \tag{Table 7, rule 1} \\
%         & \equiv \neg \neg p \vee \neg (p \vee q) \vee q \tag{DeMorgan's law} \\
%         & \equiv p \vee \neg (p \vee q) \vee q \tag{Double Negation} \\
%         & \equiv p \vee q \vee \neg (p \vee q) \tag{Commutative} \\
%         & \equiv (p \vee q) \vee \neg (p \vee q) \tag{algebra, add ()} \\
%         & \equiv \mathbf{T} \tag{Negation} 
%     \end{align*}

%     \begin{center}
%     OR
%     \end{center}

%     \vspace{-20pt}
%     \begin{align*}
%         [ \neg p & \wedge (p \vee q) ] \ra q \\
%         & \equiv \neg [ \neg p \wedge (p \vee q) ] \vee q \tag{Table 7, rule 1} \\
%         & \equiv \neg \neg p \vee \neg (p \vee q) \vee q \tag{DeMorgan's law} \\
%         & \equiv p \vee \neg (p \vee q) \vee q \tag{Double Negation} \\
%         & \equiv p \vee q \vee \neg (p \vee q) \tag{Commutative} \\
%         & \equiv p \vee q \vee (\neg p \wedge \neg q) \tag{DeMorgan's} \\
%         & \equiv p \vee (q \vee \neg p) \wedge (q \vee \neg q) \tag{Distributive} \\
%         & \equiv p \vee (q \vee \neg p) \wedge \mathbf{T} \tag{Negation} \\
%         & \equiv p \vee (q \vee \neg p) \tag{Identity} \\
%         & \equiv p \vee (\neg p \vee q) \tag{Commutative} \\
%         & \equiv (p \vee \neg p) \vee q \tag{Associative} \\
%         & \equiv \mathbf{T} \vee q \tag{Negation} \\
%         & \equiv \mathbf{T} \tag{Domination}
%     \end{align*}
% \end{solution}


% % CS 2311 f14 hw5
% \ugquestion{8} 
% Show the following statement is a tautology without using truth tables. Use the equivalence laws from Table 6-8 (model the solutions in the style of book Examples 6-8, pp. 29-30).  Justify each step with laws.

% \vspace{-0.15in}
% \begin{center}
%     $[(p \vee q) \wedge (p \ra r) \wedge (q \ra r)] \ra r$
% \end{center}
% \vspace{-10pt}

%     \begin{solution}
%     \begin{align*}
%       [(p \vee q) &\wedge ((p \ra r) \wedge (q \ra r))] \ra r \\
%       &\equiv \neg[(p \vee q) \wedge \neg(p \vee r) \wedge \neg (q \vee r)] \vee r \tag{Table 7, rule 1, x3} \\
%       &\equiv [\neg(p \vee q) \vee \neg(\neg(p \vee r) \wedge \neg(q \vee r))] \vee r \tag{DeMorgan's} \\
%       & \equiv (\neg p \wedge \neg q) \vee (\neg \neg(p \vee r) \vee \neg \neg(q \vee r)) \vee r \tag{DeMorgan's, x2} \\
%       & \equiv (\neg p \wedge \neg q) \vee (p \vee r) \vee (q \vee r) \vee r \tag{Double Negation, x2} \\
%       & \equiv (\neg p \wedge \neg q) \vee p \vee q \vee r \vee r \vee r \tag{Commutative} \\
%       & \equiv (\neg p \wedge \neg q) \vee p \vee q \vee r \vee r \tag{Idempotent} \\
%       & \equiv (\neg p \wedge \neg q) \vee p \vee q \vee r \tag{Idempotent} \\
%       & \equiv p \vee (\neg p \wedge \neg q) \vee q \vee r \tag{Commutative} \\
%       & \equiv (p \vee \neg p) \wedge (p \vee \neg q) \vee q \vee r \tag{Distributive} \\
%       & \equiv (\mathbf{T} \wedge (p \vee \neg q) \vee q \vee r \tag{Negation} \\
%       & \equiv (p \vee \neg q) \vee q \vee r \tag{Identity} \\
%       & \equiv p \vee (\neg q \vee q) \vee r \tag{Associative} \\
%       & \equiv p \vee \mathbf{T} \vee r \tag{Negation} \\
%       & \equiv \mathbf{T} \vee r \tag{Domination} \\
%       & \equiv \mathbf{T} \tag{Domination}
%     \end{align*}

%     An alternative solution is:
%     \begin{align*}
%       [(p \vee q) &\wedge ((p \ra r) \wedge (q \ra r))] \ra r \\
%        &\equiv [(p \vee q) \wedge ((p \vee q) \ra r)] \ra r \tag{Table 7, rule 7} \\
%        &\equiv [(p \vee q) \wedge (\neg (p \vee q) \vee r)] \ra r \tag{Table 7, rule 1} \\
%        &\equiv [((p \vee q) \wedge \neg (p \vee q)) \vee ((p \vee q) \wedge r)] \ra r \tag{Distributive} \\
%        &\equiv [\mathbf{F} \vee ((p \vee q) \wedge r)] \ra r \tag{Negation} \\
%        & \equiv [((p \vee q) \wedge r) \vee \mathbf{F}] \ra r \tag{Commutative} \\
%        & \equiv [((p \vee q) \wedge r)] \ra r \tag{Identity} \\
%        & \equiv [(p \vee q) \wedge r] \ra r \tag{drop extra ()} \\
%        & \equiv \neg [(p \vee q) \wedge r] \vee r \tag{Table 7, rule 1} \\
%        & \equiv  [\neg (p \vee q) \vee \neg r] \vee r \tag{DeMorgans} \\
%        & \equiv  \neg (p \vee q) \vee [\neg r \vee r ] \tag{Associative} \\
%        & \equiv \neg (p \vee q) \vee [r \vee \neg r] \tag{Commutative} \\
%        & \equiv \neg (p \vee q) \vee \mathbf{T} \tag{Negation} \\
%        & \equiv \mathbf{T} \tag{Domination} 
%     \end{align*}
%     \end{solution}

% % CS 2311 f16
% \ifprintanswers
% \newpage
% \fi
% \ugquestion{8} Prove the following statement is a tautology without using truth tables (use the logical equivalences from Table 6-8 of the book).  Justify each step with the law used.  Model the solutions in the style of Examples 6-8, pp. 29-30 of the book.
% \[ [(p \rightarrow q) \wedge (q \rightarrow r)] \rightarrow (p \rightarrow r) \]
%     \ifprintanswers
%         \vspace{-30pt}
%     \fi
% 	\begin{solution}
%    \begin{align*}
%     [(p \rightarrow q) &\wedge (q \rightarrow r)] \rightarrow (p \rightarrow r) \\
%      & \equiv \neg [ (\neg p \vee q) \wedge (\neg q \vee r)] \vee (\neg p \vee r) \tag{Table 7, rule 1, x4}\\
%      & \equiv [\neg (\neg p \vee q) \vee \neg (\neg q \vee r)] \vee (\neg p \vee r) \tag{DeMorgan's} \\
%      & \equiv [(\neg \neg p \wedge \neg q) \vee (\neg \neg q \wedge \neg r)] \vee (\neg p \vee r) \tag{DeMorgan's, x2} \\
%      & \equiv (p \wedge \neg q) \vee (q \wedge \neg r) \vee \neg p \vee r \tag{Double Negation, x2} \\
%      & \equiv (p \wedge \neg q) \vee \neg p \vee (q \wedge \neg r) \vee r \tag{Commutative} \\
%      & \equiv \neg p \vee (p \wedge \neg q) \vee r \vee (q \wedge \neg r) \tag{Commutative, x2} \\
%      & \equiv (\neg p \vee p) \wedge (\neg p \vee \neg q) \vee (r \vee q) \wedge (r \vee \neg r) \tag{Distributive, x2} \\
%      & \equiv \mathbf{T} \wedge (\neg p \vee \neg q) \vee (r \vee q) \wedge \mathbf{T} \tag{Negation, x2} \\
%      & \equiv (\neg p \vee \neg q) \vee (r \vee q) \tag{Identity, x2} \\
%      & \equiv \neg p \vee \neg q \vee r \vee q \tag{algebra} \\
%      & \equiv \neg p \vee r \vee \neg q \vee q \tag{Commutative} \\
%      & \equiv \neg p \vee r \vee \mathbf{T} \tag{Negation} \\
%      & \equiv \neg p \vee \mathbf{T} \tag{Domination} \\
%      & \equiv \mathbf{T} \tag{Domination}
%   \end{align*}


%   \begin{center}
%   OR
%   \end{center}
%   \vspace{-20pt}
%   \begin{align*}
%     [(p \rightarrow q) \wedge (q \rightarrow r)] &\rightarrow (p \rightarrow r) \\
%     & \equiv [(p \vee q) \ra r] \ra (p \ra r) \tag{Table 7, r7} \\
%     & \equiv \neg[\neg(p \vee q) \vee r] \vee (\neg p \vee r) \tag{Table 7, r1, x3} \\
%     & \equiv [\neg \neg (p \vee q) \vee r] \vee (\neg p \vee r) \tag{DeMorgans} \\
%     & \equiv [(p \vee q) \vee r] \vee (\neg p \vee r) \tag{Double Negation} \\
%     & \equiv p \vee q \vee r \vee \neg p \vee r  \tag{algebra} \\
%     & \equiv p \vee q \vee \neg p \vee r \vee r \tag{Commutative} \\
%     & \equiv p \vee \neg p \vee q \vee r \vee r \tag{Commutative} \\
%     & \equiv \mathbf{T} \vee q \vee r \vee r \tag{Negation} \\
%     & \equiv \mathbf{T} \vee q \vee r \tag{Idempotent} \\
%     & \equiv q \vee \mathbf{T} \vee r \tag{Commutative} \\
%     & \equiv \mathbf{T} \vee r \tag{Domination} \\
%     & \equiv r \vee \mathbf{T} \tag{Commutative} \\
%     & \equiv \mathbf{T} \tag{Domination}
%   \end{align*}
%   \end{solution}


% CS 2311 f16
% \ifprintanswers
% \newpage
% \fi
% \gquestion{18}{10}{a-b} Prove the following logical equivalences using the
% equivalences from Table 6-8 (model the solutions in the style of
% book Examples 6-8, pp. 29-30).  Justify each step with laws.
% \begin{enumerate}[label=(\alph*),itemsep=0pt,parsep=0pt,topsep=0pt,partopsep=0pt]
%     \item $(p \wedge \neg(q \vee \neg p)) \equiv p \wedge \neg q$
%     \item $p \rightarrow (\neg q \rightarrow r) \equiv \neg(q \vee r)
%     \rightarrow \neg p$
%     \item $\neg(\neg p \vee (p \wedge q)) \vee q \equiv q \vee p $
% \end{enumerate}

%     \begin{solution}
%     \begin{parts}
%     \part  $(p \wedge \neg(q \vee \neg p)) \equiv 
%       p \wedge \neg q$
%     \begin{align*}
%       (p &\wedge \neg(q \vee \neg p)) \\
%       & \equiv p \wedge (\neg q \wedge \neg \neg p) \tag{DeMorgan's} \\
%       & \equiv p \wedge \neg q \wedge p \tag{Double Negation} \\
%       & \equiv p \wedge p \wedge \neg q \tag{Commutative} \\
%       & \equiv p \wedge \neg q \tag{Idempotent}
%     \end{align*}
    
%     \part $p \rightarrow (\neg q \rightarrow r) \equiv 
%       \neg(q \vee r) \rightarrow \neg p$
%     \begin{align*}
%       p & \rightarrow (\neg q \rightarrow r) \\
%       & \equiv p \ra (\neg \neg q \vee r) \tag{Table 7, rule 1} \\
%       & \equiv p \ra (q \vee r) \tag{Double Negation} \\
%       & \equiv \neg (q \vee r) \ra \neg p \tag{Table 7, rule 2} \\
%     \end{align*}

%     \vspace{-30pt}
%     \begin{center}
%     OR
%     \end{center}

%     \vspace{-30pt}
%     \begin{align*}
%       p & \rightarrow (\neg q \rightarrow r) \\
%        & \equiv p \ra (q \vee r) \tag{Table 7, rule 3} \\
%        & \equiv \neg (q \vee r) \ra \neg p  \tag{Table 7, rule 2}
%     \end{align*}
    
%     \part $\neg(\neg p \vee (p \wedge q)) \vee q \equiv q \vee p $
%     \begin{align*}
%     \neg(\neg p & \vee (p \wedge q)) \vee q \\
%      & \equiv (\neg \neg p \wedge \neg (p \wedge q)) \vee q \tag{DeMorgan's} \\
%      & \equiv (p \wedge \neg (p \wedge q)) \vee q \tag{Double Negation}\\
%      & \equiv (p \wedge (\neg p \vee \neg q)) \vee q \tag{DeMorgan's}\\
%      & \equiv (p \wedge \neg p) \vee (p \wedge \neg q) \vee q \tag{Distributive}\\
%      & \equiv \mathbf{F} \vee (p \wedge \neg q) \vee q \tag{Negation}\\
%      & \equiv (p \wedge \neg q) \vee q \tag{Identity}\\
%      & \equiv q \vee (p \wedge \neg q) \tag{Commutative}\\
%      & \equiv (q \vee p) \wedge (q \vee \neg q) \tag{Distributive}\\
%      & \equiv (q \vee p) \wedge \mathbf{T} \tag{Negation}\\
%      & \equiv  q \vee p \tag{Identity} \\
%      & \equiv p \vee q \tag{Commutative}
%      \end{align*}

%     \vspace{-30pt}
%     \begin{center}
%     OR
%     \end{center}

%     \vspace{-30pt}
%     \begin{align*}
%     \neg(\neg p & \vee (p \wedge q)) \vee q \\
%       & \equiv \neg[ (\neg p \vee p) \wedge (\neg p \vee q) ] \vee q \tag{Distributive} \\
%       & \equiv \neg[ \mathbf{T} \wedge (\neg p \vee q) ] \vee q \tag{Negation} \\
%       & \equiv \neg[ (\neg p \vee q)  ] \vee q \tag{Identity} \\
%       & \equiv \neg( \neg p \vee q) \vee q \tag{algebra} \\
%       & \equiv (\neg \neg p \wedge \neg q) \vee q \tag{DeMorgans} \\
%       & \equiv (p \wedge \neg q) \vee q \tag{Double Negation} \\
%       & \equiv q \vee (p \wedge \neg q) \tag{Commutative} \\
%       & \equiv (q \vee p) \wedge (q \vee \neg q) \tag{Distributive} \\
%       & \equiv (q \vee p) \wedge \mathbf{T} \tag{Negation} \\
%       & \equiv (q \vee p) \tag{Identity} \\
%       & \equiv q \vee p \tag{algebra}
%     \end{align*}
%     \end{parts}
%     \end{solution}



% \newpage
\ugquestion{8}  Rosen Ch 1.6 \#6, p. 82.\\
Let $r$ be the proposition ``It rains", let $f$ be ``It is foggy",
    let $s$ be ``The sailing race will be held", let $l$ be ``The life
    saving demonstration will go on", and let $t$ be ``The trophy will
    be awarded".
    \ifprintanswers
        \vspace{-10pt}
    \fi
\begin{solution}
    The premises are: $(\neg r \vee \neg f) \rightarrow (s \wedge l)$, $s \rightarrow t$, $\neg t$. \\
    The conclusion we want is: $r$.

    \begin{tabular}{lll}
        Step    & \hspace{0.2in} & Reason \\
        \hline
        1. $\neg t$                 &       & Hypothesis (Premise, Given) \\
        2. $s \rightarrow t$        &       & Hypothesis (Premise, Given) \\
        3. $\neg s$                 &       & Modus tollens with (1) and (2) \\
        4. $(\neg r \vee \neg f) \rightarrow (s \wedge l)$  &   & Hypothesis (Premise, Given)  \\
        5. $(\neg(s \wedge l)) \rightarrow \neg(\neg r \vee \neg f)$    & & Contrapositive with (4) \\
        6. $(\neg s \vee \neg l) \rightarrow (\neg \neg r \wedge \neg \neg f)$ & & De Morgans law with (5) (x2) \\
        7. $(\neg s \vee \neg l) \rightarrow (r \wedge f)$  & & Double negation with (6) (x2) \\
        8. $\neg s \vee \neg l$     &       & Addition with (3) \\
        9. $r \wedge f$             &       & Modus ponens with (7) and (8) \\
        10. $r$                     &       & Simplification with (9)
    \end{tabular}

    \medskip
    Another valid argument is:

    \begin{tabular}{lll}
        Step   & \hspace{0.2in}     & Reason \\
        \hline
        1. $\neg t$                 &   & Hypothesis \\
        2. $s \ra t$                &   & Hypothesis \\
        3. $(\neg r \vee \neg f) \ra (s \wedge l)$  &   & Hypothesis \\
        4. $\neg s$                 &   & Modus tollens, (1), (2) \\
        5. $\neg (\neg r \vee \neg f) \vee (s \wedge l) $  & & Table 7, rule 1, (3) \\
        6. $(\neg \neg r \wedge \neg \neg f) \vee (s \wedge l)$  & & DeMorgans,  (5) \\
        7. $(r \wedge f) \vee (s \wedge l)$     & & Double Negation (x2), (6) \\
        8. $\neg s \vee \neg l$     &   & Addition with (4) \\
        9. $\neg (s \wedge l)$      &   & DeMorgans, (8) \\
        10. $r \wedge f$            &   & Disjunctive Syllogism, (7), (9) \\
        11. $r$                     &   & Simplification, (10)
    \end{tabular}


    \newpage
    Or, another valid argument is:

    \begin{tabular}{lll}
        Step   & \hspace{0.2in}     & Reason \\
        \hline
        1. $(\neg r \vee \neg f) \ra (s \wedge l)$                &   & Hypothesis \\
        2. $s \ra t$                &   & Hypothesis \\
        3. $\neg t$  &   & Hypothesis \\
        4. $\neg s$                 &   & Modus tollens, (2), (3) \\
        5. $\neg s \vee \neg l$     &   & Addition with (4) \\
        6. $\neg(s \wedge l)$       &   & DeMorgans with (5) \\
        7. $\neg (\neg r \vee \neg f)$  & & Modus tollens with (1) and (6)\\
        8. $\neg \neg r \wedge \neg \neg f$ & & DeMorgans with (7) \\
        9. $r \wedge f$             &   & Double Negation with (8) \\
        10. $r$                     &   & Simplification, (9)
    \end{tabular}
\end{solution}



% \gquestion{10}{10}{all} Use the rules of inference to show that the hypotheses imply the conclusion:
% \begin{itemize}[itemsep=0pt,parsep=0pt,topsep=0pt,partopsep=0pt]
%     \item ``If I graduate in four years, then I will have completed the CS courses", and
%     \item ``If I do not work on CS for 10 hours a week, then I will not complete the CS courses", and
%     \item ``If I work on CS for 10 hours a week, then I can not procrastinate."
% \end{itemize}
% Conclusion: ``If I procrastinate, then I will not graduate in four years."
% Let
% \begin{itemize}[itemsep=0pt,parsep=0pt,topsep=0pt,partopsep=0pt]
%     \item[$w = $] ``I work on CS for 10 hours a week",
%     \item[$g = $] ``I graduate in four years",
%     \item[$c = $] ``I will complete the CS courses", and
%     \item[$p = $] ``I procrastinate."
% \end{itemize}

% \begin{parts}
%   \part (4 pts) Translate the hypotheses and conclusion to logical statements. 
%   \part (6 pts) Construct a valid argument (justify each step). 
% \end{parts}

% \textit{Hint: Remember you can also use the logical equivalences as a step in the argument.}
%     \ifprintanswers
%         \vspace{-12pt}
%     \fi
% \begin{solution}
%   (a) The hypotheses are: $g \ra c$, $\neg w \ra \neg c$ and $w \ra \neg p$. \\
%   The conclusion is: $p \ra \neg g$.
  
%   (b)

%     \begin{tabular}{lll}
%         Step    & \hspace{0.2in} & Reason \\
%         \hline
%         1. $g \ra c$                & & hypothesis \\
%         2. $\neg w \ra \neg c$          & & hypothesis \\
%         3. $w \ra \neg p$             & & hypothesis \\
%         4. $\neg \neg p \ra \neg w$   & & Table 7, rule 2, with (3) \\
%         5. $p \ra \neg w$       & & Double Negation with (4) \\
%         6. $p \ra \neg c$               & & Hyp. syl., with (2) \& (5) \\
%         7. $\neg c \ra \neg g$          & & Table 7, rule 2 with (1) \\
%         8. $p \ra \neg g$       & & Hyp. syl. with (6) \& (7)
%     \end{tabular}
    
%     \emph{Note, this is one possible valid argument; many others exists.}

%     \begin{tabular}{lll}
%         Step    & \hspace{0.2in} & Reason \\
%         \hline
%         1. $g \ra c$                    & & premise \\
%         2. $\neg w \ra \neg c$          & & premise \\
%         3. $w \ra \neg p$               & & premise \\
%         4. $\neg g \vee c$              & & Table 7.1 with (1) \\
%         5. $\neg \neg w \vee \neg c$    & & Table 7.1 with (2) \\
%         6. $\neg w \vee \neg p$         & & Table 7.1 with (3) \\
%         7. $w \vee \neg c$              & & Double Negation with (5) \\
%         8. $\neg p \vee \neg c$         & & Resolution with (6) and (7) \\
%         9. $\neg c \vee \neg p$         & & Commutative with (8) \\
%         10. $c \vee \neg g$             & & Commutative with (9) \\
%         11. $\neg p \vee \neg g$        & & Resolution with (9) and (10) \\
%         12. $p \ra \neg g$              & & Table 7.1 with (11)
%     \end{tabular}
% \end{solution}


% \gquestion{12}{12}{all} Use the rules of inference to show that the hypotheses imply the conclusion:

% \begin{quote}
% If you eat carefully then you will have a healthy digestive system. If you exercise regularly you will be very fit. If you have a healthy digestive system or you are very fit, you will live to a ripe old age. You do not live to a ripe old age. Therefore, you did not eat carefully and you did not exercise regularly.
% \end{quote}

% Let
% \begin{itemize}[itemsep=0pt,parsep=0pt,topsep=0pt,partopsep=0pt]
%     \item[$c = $] ``you eat carefully",
%     \item[$h = $] ``you have a healthy digestive system",
%     \item[$e = $] ``you excercise regularly",
%     \item[$f = $] ``you will be very fit", 
%     \item[$l = $] ``you will live to a ripe old age"
% \end{itemize}

% \begin{parts}
%   \part (4 pts) Translate the hypotheses and conclusion to logical statements. 
%   \part (8 pts) Construct a valid argument (justify each step). 
% \end{parts}

% \textit{Hint: Remember you can also use the logical equivalences as a step in the argument.}
%     \ifprintanswers
%         \vspace{-2pt}
%     \fi
% \begin{solution}
%   % (a) The hypotheses are: 
%   % 	\begin{itemize}[itemsep=0pt,parsep=0pt,topsep=0pt,partopsep=0pt]
%   % 		\item $c \ra h$ 
%   % 		\item $e \ra f$ 
%   % 		\item $(h \vee f) \ra l$ 
%   % 		\item $\neg l$ 
%   % 	\end{itemize}
%   % The conclusion is: $\neg c \wedge \neg e$.
%   (a) The hypotheses are:  $\;\;c \ra h$ \qquad $e \ra f$ \qquad $(h \vee f) \ra l$ 
%   \qquad $\neg l$ 

%   The conclusion is: $\neg c \wedge \neg e$
  
%   (b) 

%   \begin{tabular}{lll}
%         Step    & \hspace{0.2in} & Reason \\
%         \hline
%         1. $c \ra h$                & & hypothesis \\
%         2. $e \ra f$           & & hypothesis \\
%         3. $(h \vee f) \ra l$            & & hypothesis \\
%         4. $\neg l$ 				& & hypothesis \\
%         5. $\neg (h \vee f)$   & & modus tollens (3) \& (4) \\
%         6. $\neg h \wedge \neg f$       & & DeMorgans (5) \\
%         7. $\neg h$               & & Simplification (6) \\
%         8. $\neg c$          & & modus tollens (7) \& (1) \\
%         9. $\neg f$       & & Simplification (6) \\
%         10. $\neg e$  		& & modus tollens (2) \& (9) \\
%         11. $\neg c \wedge \neg e$ 		 & & Conjunction (8) \& (10) \\
%     \end{tabular}
    
%     \emph{Note, this is one possible valid argument; many others exists.}

%     % \begin{tabular}{lll}
%     %     Step    & \hspace{0.2in} & Reason \\
%     %     \hline
%     %     1. $c \ra h$                & & premise \\
%     %     2. $e \ra f$                & & premise \\
%     %     3. $(h \vee f) \ra l$       & & premise \\
%     %     4. $\neg l$                 & & premise \\
%     %     5. $\neg (h \vee f)$        & & modus tollens (3) \& (4) \\
%     %     6. $
% \end{solution}



\bonusquestion[2] Give a compound proposition with three variables $p$, $q$, and $r$ that is 
true when at most one of the three variables is true, and false otherwise.
    \begin{solution}
    \[ (p \wedge \neg q \wedge \neg r) \vee (\neg p \wedge q \neg r) 
        \vee (\neg p \wedge \neg q \wedge r) \vee (\neg p \wedge \neg q \neg r)\]
    \end{solution}





\end{questions}

\end{document}
